\documentclass[aspectratio=169,11pt]{beamer}
\usetheme{Madrid}
\usepackage[T1]{fontenc}
\usepackage[utf8]{inputenc}
\usepackage[french]{babel}
\usepackage{amsmath,amssymb,stmaryrd}
\usepackage{tikz}
\usepackage{fontawesome5}
\usepackage{booktabs}
\usepackage{listings}

\lstset{basicstyle=\ttfamily\footnotesize, columns=fullflexible, keepspaces=true,
        keywordstyle=\bfseries, language=Python}

\setbeamertemplate{section in toc}[sections numbered]
\title{Un solveur SAT pour Meowdoku}
\subtitle{MLSF -- Projet, partie I}
\author{Salomé Fonvielle \and Hugo Vigna}
\date{Octobre 2026}

% --- Grilles -------------------------------------------------------------
\definecolor{regA}{HTML}{F4CCCC}
\definecolor{regB}{HTML}{CFE2F3}
\definecolor{regC}{HTML}{D9EAD3}
\definecolor{regD}{HTML}{FFF2CC}
\definecolor{regE}{HTML}{D9D2E9}
\definecolor{regF}{HTML}{FCE5CD}

% \grille[échelle]{lignes}{chats} : lignes = {{A,B,...},...}, chats = liste i/j (1-indexés), vide = aucun
\newcommand{\grille}[3][0.5]{%
  \begin{tikzpicture}[scale=#1]
    \foreach \ligne [count=\i] in #2 {
      \foreach \z [count=\j] in \ligne {
        \fill[reg\z] (\j-1,-\i) rectangle ++(1,1);
        \node[font=\tiny, text=black!50] at (\j-0.75,-\i+0.75) {\z};
      }
    }
    \draw[black!25] (0,-6) grid (6,0);
    \draw[thick] (0,-6) rectangle (6,0);
    \foreach \i/\j in {#3} {
      \if\relax\detokenize\expandafter{\i}\relax\else
        \node at (\j-0.5,-\i+0.5) {\faCat};
      \fi
    }
  \end{tikzpicture}}

\newcommand{\sujet}{{A,B,C,B,B,D},{A,B,B,B,B,D},{A,A,B,D,D,D},{A,A,A,E,D,D},{A,A,A,A,D,D},{F,A,D,D,D,D}}
\newcommand{\genere}{{B,A,A,A,C,C},{B,B,A,C,C,C},{B,B,C,C,C,C},{E,D,D,D,D,C},{E,E,D,F,F,C},{E,E,F,F,F,F}}

\begin{document}

\maketitle

\begin{frame}{Plan}
  \tableofcontents
\end{frame}

% =========================================================================
\section{Modéliser Meowdoku en CNF}

\begin{frame}{Le jeu et ce qui est demandé}
  \begin{columns}[T]
    \column{0.52\textwidth}
    \begin{block}{Règles}
      Grille $N\times N$, $N$ régions, $N$ chats
      {\footnotesize (on se restreint à $N = M = K$)} :
      \begin{itemize}
        \item au plus un chat par ligne et par colonne ;
        \item exactement un chat par région ;
        \item aucun chat collé à un autre (8-voisinage).
      \end{itemize}
    \end{block}
    \begin{block}{Demandé}
      \begin{enumerate}
        \item Grille $\to$ DIMACS (variables $x_{i,j}$) dont toute valuation témoin
          est une solution.
        \item Générateur de grilles à solution garantie.
      \end{enumerate}
      {\footnotesize Fichiers : \texttt{meowdoku.py}, \texttt{generator.py}}
    \end{block}
    \column{0.3\textwidth}
    \centering
    \grille[0.4]{\sujet}{} \\ {\footnotesize grille du sujet} \\[4pt]
    \grille[0.4]{\sujet}{3/2,2/5,1/3,5/6,4/4,6/1} \\ {\footnotesize une solution}
  \end{columns}
\end{frame}

\begin{frame}[fragile]{Variables}
  \begin{columns}[T]
    \column{0.6\textwidth}
    \begin{block}{Une variable par case}
      $x_{i,j}$ vraie ssi un chat est en $(i,j)$. DIMACS numérote à partir de 1 :
      on cherche une bijection
      \[ \varphi : \llbracket 1,N \rrbracket^2 \to \llbracket 1,N^2 \rrbracket \]
      \[ \varphi(i,j) = (i-1)\,N + j \]
      \[ \varphi^{-1}(k) = \big(\lfloor \tfrac{k-1}{N} \rfloor + 1,\; (k-1) \bmod N + 1\big) \]
    \end{block}
    \begin{block}{Dans le code (indices à partir de 0)}
\begin{lstlisting}
def var(i, j, n):        # = phi(i+1, j+1)
    return i * n + j + 1
\end{lstlisting}
    \end{block}
    \column{0.32\textwidth}
    \centering
    \begin{tikzpicture}[scale=0.55]
      \foreach \i in {1,...,6}
        \foreach \j in {1,...,6} {
          \pgfmathtruncatemacro{\k}{(\i-1)*6+\j}
          \node[font=\footnotesize] at (\j-0.5,-\i+0.5) {\k};
        }
      \draw[black!25] (0,-6) grid (6,0);
      \draw[thick] (0,-6) rectangle (6,0);
      \foreach \j in {1,...,6} \node[font=\tiny, text=black!50] at (\j-0.5,0.4) {$j{=}\j$};
      \foreach \i in {1,...,6} \node[font=\tiny, text=black!50, anchor=east] at (0,-\i+0.5) {$i{=}\i$};
    \end{tikzpicture} \\[2pt]
    \footnotesize numérotation pour $N = 6$
  \end{columns}
\end{frame}

\begin{frame}{Clauses et correction}
  \begin{enumerate}
    \item \textbf{Conflit} : $\neg x_a \lor \neg x_b$ pour toute paire $\{a,b\}$ sur la même ligne, la même colonne,
      dans la même région, ou adjacente (8-voisinage).
    \item \textbf{Région} : $\bigvee_{a \in R_k} x_a$ pour toute région $R_k$.
  \end{enumerate}
  \medskip
  \begin{center}\small
  \begin{tabular}{ll}
    \toprule
    Règle du sujet & Clauses \\
    \midrule
    au plus un chat par ligne & conflit « même ligne » \\
    au plus un chat par colonne & conflit « même colonne » \\
    exactement un chat par région & région + conflit « même région » \\
    aucun chat adjacent & conflit « adjacentes » \\
    \bottomrule
  \end{tabular}
  \end{center}
  Donc \textbf{modèles de la CNF $=$ configurations valides}.
  \medskip

  {\footnotesize « Au plus un » par paires : $\frac{n(n-1)}{2}$ clauses, sans variable auxiliaire ;
  poser $a$ rend chaque $\neg a \lor \neg b$ unitaire (propagation).
  Avec $N = M = K$, « exactement un par ligne/colonne » en découle.\par}
\end{frame}

\begin{frame}{Exemple : la grille du sujet}
  \begin{columns}[c]
    \column{0.36\textwidth}
    \centering
    \begin{tikzpicture}[scale=0.62]
      \foreach \ligne [count=\i] in \sujet {
        \foreach \z [count=\j] in \ligne {
          \fill[reg\z] (\j-1,-\i) rectangle ++(1,1);
          \pgfmathtruncatemacro{\k}{(\i-1)*6+\j}
          \node[font=\footnotesize] at (\j-0.5,-\i+0.5) {\k};
        }
      }
      \draw[black!25] (0,-6) grid (6,0);
      \draw[thick] (0,-6) rectangle (6,0);
      \foreach \i/\j in {1/3,4/4,6/1} \draw[red, very thick] (\j-1,-\i) rectangle ++(1,1);
    \end{tikzpicture} \\[2pt]
    \footnotesize numéro de variable de chaque case
    \column{0.6\textwidth}
    {\small\begin{tabular}{@{}l@{\quad}l@{}}
      \texttt{p cnf 36 309} & 36 variables, 309 clauses \\[4pt]
      \texttt{-1 -2 0}      & pas de chat en 1 \emph{ou} pas en 2 (ligne) \\
      \texttt{-1 -7 0}      & pas de chat en 1 \emph{ou} pas en 7 (colonne) \\
      \texttt{...}          & 303 clauses de conflit \\[4pt]
      \texttt{1 7 13 ... 32 0} & au moins un chat dans A \\
      \texttt{3 0}          & C n'a qu'une case : chat en 3 \\
      \texttt{22 0}         & E : chat en 22 \\
      \texttt{31 0}         & F : chat en 31 \\
      \texttt{...}          & 6 clauses de région \\
    \end{tabular}}
    \bigskip

    Les cases \textcolor{red}{encadrées} sont des régions d'une seule case :
    leur clause est unitaire, DPLL y place un chat sans décision.
  \end{columns}
\end{frame}

\begin{frame}{Générateur : solution plantée}
  \begin{columns}[c]
    \column{0.62\textwidth}
    \begin{enumerate}
      \item \textbf{Placer les chats} : une colonne par ligne, toutes distinctes,
        avec un écart $\ge 2$ entre lignes consécutives (backtracking aléatoire).
        Impossible pour $N = 2, 3$.
      \item \textbf{Faire grossir les régions} : chaque chat est la graine d'une région.
        Tant qu'il reste une case libre, on en tire une au hasard parmi celles qui touchent
        une région (4-voisinage) et on l'y rattache.
    \end{enumerate}
    \medskip
    Les régions obtenues sont connexes et contiennent chacune exactement un chat.
    La configuration plantée est donc valide.
    \column{0.3\textwidth}
    \centering \grille{\genere}{1/3,2/1,3/6,4/4,5/2,6/5} \\
    \footnotesize $N=6$, graine 1
  \end{columns}
\end{frame}

\begin{frame}{Validation}
  \begin{itemize}
    \item \texttt{is\_valid} vérifie les règles \textbf{directement}, sans passer par la CNF.
    \item Grilles générées pour $N = 4..8$, 20 graines chacune : la solution plantée est valide,
      elle satisfait la CNF, et les régions sont connexes.
    \item \textbf{Équivalence exhaustive} pour $N = 4,5,6$ : pour chacune des $N!$ permutations,
      « satisfait la CNF » $\Leftrightarrow$ « \texttt{is\_valid} ».
    \item Aller-retour fichier : \texttt{write\_grid} puis \texttt{parse\_grid} redonne la grille.
  \end{itemize}
\end{frame}

\begin{frame}{Choix et limites}
  \begin{itemize}
    \item $N = M = K$ uniquement.
    \item La solution n'est \textbf{pas unique} : 20 solutions pour la grille $N=6$, graine 1.
      Le sujet n'exige qu'une solution au moins.
    \item Encodage par paires : $O(|R|^2)$ clauses par région ; une grande région coûte cher.
    \item Piste pour la Q6 : ajouter les clauses redondantes « au moins un par ligne/colonne »
      pour obtenir plus de propagation unitaire.
  \end{itemize}
\end{frame}

% =========================================================================
\section{DIMACS et DPLL}
\begin{frame}{À compléter}\end{frame}

\section{Résoudre une grille}
\begin{frame}{À compléter}\end{frame}

\section{Heuristiques}
\begin{frame}{À compléter}\end{frame}

\section{Comparaison}
\begin{frame}{À compléter}\end{frame}

\section{Solveur optimisé}
\begin{frame}{À compléter}\end{frame}

\section{Usage de l'IA}
\begin{frame}{À compléter}\end{frame}

\end{document}
